%%%%%%%%%%%%%%%%%%%%%%%%%%%%%%%%%%%%%%%%%%%%%%%%%%%%%%%%%%%%%%%%%%%%%%%%%%%
%
% Generic template for TFC/TFM/TFG/Tesis
%
% $Id: manual.tex,v 1.14 2016/03/31 10:44:12 macias Exp $
%
% By:
%  + Javier Macías-Guarasa. 
%    Departamento de Electrónica
%    Universidad de Alcalá
%  + Roberto Barra-Chicote. 
%    Departamento de Ingeniería Electrónica
%    Universidad Politécnica de Madrid   
% 
% Based on original sources by Roberto Barra, Manuel Ocaña, Jesús Nuevo,
% Pedro Revenga, Fernando Herránz and Noelia Hernández. Thanks a lot to
% all of them, and to the many anonymous contributors found (thanks to
% google) that provided help in setting all this up.
%
% See also the additionalContributors.txt file to check the name of
% additional contributors to this work.
%
% If you think you can add pieces of relevant/useful examples,
% improvements, please contact us at (macias@depeca.uah.es)
%
% You can freely use this template and please contribute with
% comments or suggestions!!!
%
%%%%%%%%%%%%%%%%%%%%%%%%%%%%%%%%%%%%%%%%%%%%%%%%%%%%%%%%%%%%%%%%%%%%%%%%%%%


\chapter{Código desarrollado}
\label{code}

Todo el código desarrollado para este proyecto, tanto la implementación de las metaheurísticas expuesta en el
capítulo~\ref{cha:design}, como los programas de Python auxiliares para el análisis comparativo y generación de gráficas,
se puede encontrar en el siguiente repositorio de GitHub:

\url{https://github.com/ARetamosaG/TFG}

En este repositorio se puede encontrar lo siguiente:

\begin{itemize}
  \item Un subdirectorio \textit{algorithms}: contiene el código fuente de las cuatro metaheurísticas.
  \item Un subdirectorio \textit{experiment\_results}: contiene los ficheros \textit{csv} auxiliares generados durante el 
  experimento. Como apunte, el fichero \textit{history.csv} es considerablemente extenso, puesto que contiene una línea por 
  cada iteración/generación habida durante todo el experimento. Al ejecutar seis instancias, cinco veces cada una, con cuatro
  metaheurísticas, en total hay 120 ejecuciones, cada una con decenas o cientos de iteraciones.
  \item Un subdirectorio \textit{graphs}: contiene todas las gráficas y diagramas mostradas en este documento.
  \item Un subdirectorio \textit{instances}: contiene las instancias de TSPLIB empleadas en el experimento. No obstante, estas 
  se pueden obtener desde \cite{tsplib}, aunque algunas se extrajeron desde \cite{shredder_github}, debido a que no estaban
  disponibles en la página oficial.
  \item Un subdirectorio \textit{scripts}: contiene una serie de ficheros \textit{.py}, encargados de ejecutar el experimento,
  analizar los ficheros \textit{csv} resultado, generar las gráficas y calcular el tiempo medio de ejecución por instancia y 
  algoritmo.
\end{itemize}

Asimismo, también se pueden encontrar dos ficheros \textit{.pdf} relacionados con este trabajo:
\begin{itemize}
  \item El fichero de \textbf{anteproyecto}.
  \item La \textbf{memoria} del trabajo, esto es, el presente documento.
\end{itemize}

%%% Local Variables:
%%% TeX-master: "../book"
%%% End:


